\section{在线残差与驻点跟踪的证明}\label{app:online}

\subsection{归一化投影上升}

\begin{lemma}[单步归一化投影上升]\label{lem:projected-ascent}
设$J$在非空闭凸集$\Theta$上为$L$-光滑。固定$x\in\Theta$，令$g=\nabla J(x)$且$\norm{g}_2\le B_g$，$\widehat g$为任意有限向量。记
\[
 x^+=\Pi_\Theta\{x+\gamma\mathsf d(\widehat g)\},
 \qquad
 Q(x)=\gamma^{-1}[\Pi_\Theta\{x+\gamma\mathsf d(g)\}-x].
\]
若$\gamma L\le\vartheta$，且$C_g:=\max\{B_g,\vartheta\}$，则
\begin{equation}\label{eq:one-step-ascent}
 J(x^+)-J(x)
 \ge\frac{\vartheta\gamma}{8}\norm{Q(x)}_2^2
 -\frac{5C_g\gamma}{4\vartheta^2}\norm{\widehat g-g}_2^2.
\end{equation}
\end{lemma}

\begin{proof}
先用估计梯度产生的实际位移建立上升界，再将其与总体梯度产生的投影位移比较。令
\[
 \Delta x=x^+-x,
 \qquad
 c_g=\max\{\norm{g}_2,\vartheta\},
 \qquad
 \mathfrak e=\mathsf d(\widehat g)-\mathsf d(g).
\]
以可行点$x$为比较对象，投影变分不等式给出
\begin{equation}\label{eq:projection-vi}
 \langle\mathsf d(\widehat g),\Delta x\rangle\ge\frac{\norm{\Delta x}_2^2}{\gamma}.
\end{equation}
由$g=c_g\mathsf d(g)$、$L$-光滑函数的下二次界及式~\eqref{eq:projection-vi}，
\begin{align}
 J(x+\Delta x)-J(x)
 &\ge \langle g,\Delta x\rangle-\frac L2\norm{\Delta x}_2^2\notag\\
 &=c_g\langle\mathsf d(\widehat g)-\mathfrak e,\Delta x\rangle
 -\frac L2\norm{\Delta x}_2^2\notag\\
 &\ge\left(\frac{c_g}{\gamma}-\frac L2\right)\norm{\Delta x}_2^2
 -c_g\norm{\mathfrak e}_2\norm{\Delta x}_2\notag\\
 &\ge\left(\frac{3c_g}{4\gamma}-\frac L2\right)\norm{\Delta x}_2^2
 -c_g\gamma\norm{\mathfrak e}_2^2\notag\\
 &\ge\frac{c_g}{4\gamma}\norm{\Delta x}_2^2
 -c_g\gamma\norm{\mathfrak e}_2^2.
 \label{eq:actual-displacement-ascent}
\end{align}
倒数第二步使用$uv\le u^2/(4\gamma)+\gamma v^2$，最后一步使用$\gamma L\le\vartheta\le c_g$。

再令
\[
 t=\Pi_\Theta\{x+\gamma\mathsf d(g)\}-x=\gamma Q(x).
\]
欧氏投影非扩张，故$\norm{\Delta x-t}_2\le\gamma\norm{\mathfrak e}_2$。由$\norm{t}_2^2\le2\norm{\Delta x}_2^2+2\norm{\Delta x-t}_2^2$，
\[
 \norm{\Delta x}_2^2\ge\frac12\norm{t}_2^2-\gamma^2\norm{\mathfrak e}_2^2.
\]
代入式~\eqref{eq:actual-displacement-ascent}得到
\[
 J(x^+)-J(x)
 \ge\frac{c_g\gamma}{8}\norm{Q(x)}_2^2
 -\frac{5c_g\gamma}{4}\norm{\mathfrak e}_2^2.
\]
最后，$\mathsf d(g)=\Pi_{\{u:\norm{u}_2\le1\}}(g/\vartheta)$，由单位球投影非扩张性，
\[
 \norm{\mathfrak e}_2\le \vartheta^{-1}\norm{\widehat g-g}_2.
\]
对正项使用$c_g\ge\vartheta$，负项使用$c_g\le C_g=\max(B_g,\vartheta)$，即得式~\eqref{eq:one-step-ascent}。
\end{proof}

\subsection{期望平均平方投影残差}

\begin{proof}
先连接最终估计量与真实梯度。在给定$\cF_k$后，定理~\ref{thm:estimator}给出
\[
 \E(\widehat g_k^{\db}-\widetilde g_k(\theta_k)\mid\cF_k)=0,
 \qquad
 \E(\norm{\widehat g_k^{\db}-\widetilde g_k(\theta_k)}_2^2\mid\cF_k)
 \le\frac{C_{\rm var}(n,p_{\rm act},b)}{M_k}.
\]
与条件确定向量$\widetilde g_k(\theta_k)-g_k(\theta_k)$的交叉项因此为零。于是
\begin{align*}
 \E\norm{\widehat g_k^{\db}-g_k(\theta_k)}_2^2
 &=\E\norm{\widehat g_k^{\db}-\widetilde g_k(\theta_k)}_2^2
 +\E\norm{\widetilde g_k(\theta_k)-g_k(\theta_k)}_2^2\\
 &\le\frac{C_{\rm var}(n,p_{\rm act},b)}{M_k}+\E\beta_k^2,
\end{align*}
即式~\eqref{eq:true-gradient-mse}。

现在对真实时期目标$J_k$，在$x=\theta_k$和$\widehat g=\widehat g_k^{\db}$处应用引理~\ref{lem:projected-ascent}。记
\[
 \mathfrak j_k=J_k(\theta_{k+1})-J_k(\theta_k).
\]
取期望后有
\begin{equation}\label{eq:online-one-step}
 \frac{\vartheta\gamma}{8}\E\zeta_k
 \le\E \mathfrak j_k
 +\frac{5C_g\gamma}{4\vartheta^2}
 \E\norm{\widehat g_k^{\db}-g_k(\theta_k)}_2^2.
\end{equation}
跨越移动目标时，目标增量满足逐路径的望远镜恒等式
\begin{align}\label{eq:moving-objective-telescope}
 \sum_{k=1}^{K}\mathfrak j_k
 ={}&J_K(\theta_{K+1})-J_1(\theta_1)\notag\\
 &+\sum_{k=1}^{K-1}
 \{J_k(\theta_{k+1})-J_{k+1}(\theta_{k+1})\}.
\end{align}
由于$\abs{J_k}\le B_v$，第二行的上界为$\cV_K$，所以
\begin{equation}\label{eq:telescope-upper}
 \sum_{k=1}^{K}\mathfrak j_k\le2B_v+\cV_K.
\end{equation}
将式~\eqref{eq:online-one-step}逐期相加，代入式~\eqref{eq:true-gradient-mse}和式~\eqref{eq:telescope-upper}，再除以$\vartheta\gamma K/8$，得到式~\eqref{eq:online-residual-bound}。
\end{proof}

\subsection{动态局部遗憾}

\begin{proof}
按约定，$t\le0$时$Q_t(\theta_t)=0$。系数$\rho^j/A_{\rho,W}$对$j=0,\ldots,W-1$非负，且和为一。利用$x\mapsto\norm{x}_2^2$的凸性，
\[
 \norm{\bar Q_k}_2^2
 \le \frac{1}{A_{\rho,W}}\sum_{j=0}^{W-1}\rho^j
       \norm{Q_{k-j}(\theta_{k-j})}_2^2.
\]
对$k\le K$求和并重排下标，
\begin{align*}
 \Reg_{\rho,W}(K)
 &\le
 \sum_{t=1}^{K}\zeta_t
 \left\{\frac{1}{A_{\rho,W}}\sum_{j=0}^{\min\{W-1,K-t\}}\rho^j\right\}
 \le \sum_{t=1}^{K}\zeta_t.
\end{align*}
这证明式~\eqref{eq:dynamic-local-regret-bound}；取期望并除以$K$即得式~\eqref{eq:average-dynamic-local-regret-bound}。
\end{proof}

\par\smallskip\noindent\textbf{独立评估流。}\quad
对固定迭代点$\theta$，写
\[
 \mathcal Q_\theta(g)
 :=\gamma^{-1}\!\left[\Pi_\Theta\{\theta+\gamma\mathsf d(g)\}-\theta\right].
\]
由投影非扩张性与$\mathsf d$的$\vartheta^{-1}$-Lipschitz性，
\begin{equation}\label{eq:residual-evaluation-lipschitz}
 \norm{\mathcal Q_\theta(\widehat g)-\mathcal Q_\theta(g)}_2
 \le \vartheta^{-1}\norm{\widehat g-g}_2.
\end{equation}
若$\widehat g^{(1)},\widehat g^{(2)}$为相互独立的梯度估计，且满足$\E\norm{\widehat g^{(j)}-g}_2^2\le\varepsilon_g^2$，记$\widehat Q^{(j)}=\mathcal Q_\theta(\widehat g^{(j)})$，则
\begin{equation}\label{eq:cross-stream-residual-error}
 \left|
 \E\ip{\widehat Q^{(1)}}{\widehat Q^{(2)}}
 -\norm{\mathcal Q_\theta(g)}_2^2
 \right|
 \le \frac{2\varepsilon_g}{\vartheta}
 +\frac{\varepsilon_g^2}{\vartheta^2}.
\end{equation}
令$\operatorname{Bias}_j:=\E\widehat Q^{(j)}-\mathcal Q_\theta(g)$。由独立性，
\[
 \E\ip{\widehat Q^{(1)}}{\widehat Q^{(2)}}
 =\ip{\mathcal Q_\theta(g)+\operatorname{Bias}_1}
 {\mathcal Q_\theta(g)+\operatorname{Bias}_2}.
\]
Jensen不等式和式~\eqref{eq:residual-evaluation-lipschitz}给出$\norm{\operatorname{Bias}_j}_2\le\varepsilon_g/\vartheta$。又因$\norm{\mathsf d(g)}_2\le1$，有$\norm{\mathcal Q_\theta(g)}_2\le1$。故
\[
 \left|
 \ip{\mathcal Q_\theta(g)+\operatorname{Bias}_1}
 {\mathcal Q_\theta(g)+\operatorname{Bias}_2}
 -\norm{\mathcal Q_\theta(g)}_2^2
 \right|
 \le \norm{\mathcal Q_\theta(g)}_2
 (\norm{\operatorname{Bias}_1}_2+\norm{\operatorname{Bias}_2}_2)
 +\norm{\operatorname{Bias}_1}_2\norm{\operatorname{Bias}_2}_2,
\]
从而得到式~\eqref{eq:cross-stream-residual-error}。

对截至第$K$期的在线路径，令$\mathcal C_K$为独立评估抽样前由已实现迭代点和各期目标生成的$\sigma$-域。设两条完整评估流在给定$\mathcal C_K$后条件独立。对$1\le t\le K$，令$\widehat g_t^{(j)}$为第$j=1,2$条评估流在时期$t$的独立梯度估计，并定义
\[
 \widehat Q_t^{(j)}:=\mathcal Q_{\theta_t}(\widehat g_t^{(j)}),
 \qquad
 \E\!\left[\norm{\widehat g_t^{(j)}-g_t(\theta_t)}_2^2
 \mid\mathcal C_K\right]\le\varepsilon_{g,t}^2,
\]
其中$\varepsilon_{g,t}$是非负的$\mathcal C_K$-可测误差上界。对$t\le0$，令$\widehat Q_t^{(j)}=0$、$\varepsilon_{g,t}=0$，与式~\eqref{eq:dynamic-window-residual}一致。对$1\le k\le K$，定义
\[
 \widehat{\bar Q}_k^{(j)}
 =A_{\rho,W}^{-1}\sum_{r=0}^{W-1}\rho^r\widehat Q_{k-r}^{(j)},
 \qquad
 \bar\varepsilon_k
 =A_{\rho,W}^{-1}\sum_{r=0}^{W-1}\rho^r\varepsilon_{g,k-r}.
\]
条件Jensen不等式与式~\eqref{eq:residual-evaluation-lipschitz}将各平滑评估流的平均误差控制在$\bar\varepsilon_k/\vartheta$。由于$\norm{\bar Q_k}_2\le1$，结合条件独立性与相同的交叉乘积论证，
\begin{equation}\label{eq:cross-stream-dlr-error}
 \left|
 \E\!\left[\ip{\widehat{\bar Q}_k^{(1)}}{\widehat{\bar Q}_k^{(2)}}
 \mid\mathcal C_K\right]
 -\norm{\bar Q_k}_2^2
 \right|
 \le \frac{2\bar\varepsilon_k}{\vartheta}
 +\frac{\bar\varepsilon_k^2}{\vartheta^2}.
\end{equation}
再取期望，仍以上式右侧的期望控制交叉乘积的期望与$\E\norm{\bar Q_k}_2^2$的差。

\subsection{驻点集合的推论}

\begin{lemma}[驻点集合的漂移]\label{lem:stationary-set-drift}
若式~\eqref{eq:error-bound}成立，定义
\begin{equation}\label{eq:gradient-drift}
 \delta_k^g:=\sup_{\theta\in\Theta}
 \norm{g_{k+1}(\theta)-g_k(\theta)}_2.
\end{equation}
以$\dH$表示Hausdorff距离，则
\begin{equation}\label{eq:stationary-set-drift}
 \dH(\Sta_{k+1},\Sta_k)
 \le\frac{\kappa}{\vartheta}\delta_k^g.
\end{equation}
\end{lemma}

\begin{proof}[推论~\ref{cor:stationary-tracking}的证明]
由式~\eqref{eq:error-bound}，逐点有
\[
 \dist^2(\theta_k,\Sta_k)
 \le\kappa^2\norm{Q_k(\theta_k)}_2^2
 =\kappa^2\zeta_k.
\]
取期望并对$k$求平均，即得式~\eqref{eq:stationary-tracking}。进一步代入式~\eqref{eq:online-residual-bound}便得到有限期控制。
\end{proof}

\begin{proof}[引理~\ref{lem:stationary-set-drift}的证明]
归一化方向映射是将$g/\vartheta$投影至闭单位球，故具有$\vartheta^{-1}$-Lipschitz性。再由$\Pi_\Theta$非扩张和式~\eqref{eq:projected-residual}，
\begin{align*}
 \sup_{\theta\in\Theta}
 \norm{Q_{k+1}(\theta)-Q_k(\theta)}_2
 &\le
 \sup_{\theta\in\Theta}
 \norm{\mathsf d(g_{k+1}(\theta))-\mathsf d(g_k(\theta))}_2\\
 &\le\frac{\delta_k^g}{\vartheta}.
\end{align*}
若$\theta\in\Sta_k$，则$Q_k(\theta)=0$。对第$k+1$期应用式~\eqref{eq:error-bound}，得到
\[
 \dist(\theta,\Sta_{k+1})
 \le\kappa\norm{Q_{k+1}(\theta)}_2
 \le\frac{\kappa\delta_k^g}{\vartheta}.
\]
对$\theta\in\Sta_k$取上确界，得到单向Hausdorff界。交换$k$与$k+1$，由于统一梯度差异界对称，得到反向界；取两者较大值，即是式~\eqref{eq:stationary-set-drift}。
\end{proof}

\begin{proof}[推论~\ref{cor:vanishing-tracking}的证明]
在式~\eqref{eq:online-residual-bound}中，端点项为$O(K^{-1})$，而$\E\cV_K=o(K)$使目标变化项在除以$K$后趋于零。由于$C_{\rm var}(n,p_{\rm act},b)$固定，假设中的Ces\`aro条件给出
\[
 \frac1K\sum_{k=1}^{K}\left(\frac{C_{\rm var}(n,p_{\rm act},b)}{M_k}+\E\beta_k^2\right)\to0.
\]
故$\cE_K\to0$。由式~\eqref{eq:average-dynamic-local-regret-bound}得到动态局部遗憾结论；若局部误差界~\eqref{eq:error-bound}成立，再由式~\eqref{eq:stationary-tracking}得到驻点集合的结论。
\end{proof}

\par\smallskip\noindent\textbf{半合成环境中的目标变化。}\quad
令$c_k^{\rm ret}$收集半合成收益规律在第$k$期的系数。E3中期初账户状态的变化满足
\[
 \abs{\log X_{t_{k+1}}-\log X_{t_k}}
 +\abs{r_{k+1}-r_k}
 +\norm{\omega_{t_{k+1}-1}-\omega_{t_k-1}}_1
 =O(\varpi_k),
 \qquad \varpi_k=k^{-3/4},
\]
而$c_k^{\rm ret}$只在市场阶段边界及有限的过渡区间内变化，因此$\sum_k\norm{c_{k+1}^{\rm ret}-c_k^{\rm ret}}<\infty$。若冻结目标在紧分析域上对这些期初输入一致Lipschitz，则存在有限常数$L_{\rm acct},L_{\rm ret}$，使
\[
\begin{aligned}
 \sup_{\theta\in\Theta}
 \abs{J_{k+1}(\theta,r_{k+1})-J_k(\theta,r_k)}
 &\le
 L_{\rm acct}\!\left(
 \abs{\log X_{t_{k+1}}-\log X_{t_k}}
 +\abs{r_{k+1}-r_k}
 \right.\\
 &\qquad\left.
 +\norm{\omega_{t_{k+1}-1}-\omega_{t_k-1}}_1
 \right)
 +L_{\rm ret}\norm{c_{k+1}^{\rm ret}-c_k^{\rm ret}}.
\end{aligned}
\]
在这一正则性条件下，$\cV_K=O(\sum_{k<K}\varpi_k)+O(1)=O(K^{1/4})=o(K)$。该半合成构造的训练与执行使用同一条件因子块和收益规律，因此总体层面的模拟梯度差异项为零。
